\subsection{A source-level closure for the physical visibility phase}
\label{sec:physical_visibility_closure}

The two-zone inversion in Section~\ref{sec:spectral_visibility} determines a
dense share of processed ionizing photons for each stacked continuum subtype,
but it does not explain that sequence. We therefore close the observation
layer at the level of individual synthetic sources. The forward continuum is
the sum of an attenuated stellar host, a leaked accretion component, and the
cool nuclear reprocessor. We define the predicted nuclear-to-host dominance
at 5100~\AA{} as
\begin{equation}
  D \equiv
  \frac{L_{\nu,\mathrm{thermal}}(5100\,\text{\AA})}
       {L_{\nu,\mathrm{host}}(5100\,\text{\AA})}.
\end{equation}
Dense nuclear emission requires both nuclear photon dominance and
interception by the nuclear structure. To leading order, their product gives
the quadratic photon-competition closure
\begin{equation}
  \frac{Q_{\mathrm{dense}}}{Q_{\mathrm{diffuse}}}=D^2,
  \qquad
  f_{\mathrm{dense}}=\frac{D^2}{1+D^2}.
  \label{eq:visibility_closure}
\end{equation}
Equation~\ref{eq:visibility_closure} is applied before the predicted
$L_{5100}/L_{2500}$ ratio is classified. It contains no subtype-specific
normalization or dense fraction.

For the $4.5\leq z\leq6.5$ forward sample, the median predicted
$f_{\mathrm{dense}}$ values are 0.985, 0.904, 0.711, and 0.212 for xLRD,
plusLRD, minusLRD, and bLRD sources. The corresponding stack inversions are
0.981, 0.901, 0.738, and 0.086; the bLRD interval is weakly identified and
extends from zero to 0.658. Generalizing Equation~\ref{eq:visibility_closure}
to $D^\eta$ gives $\eta=1.89$ under measurement errors alone and
$\eta=2.015$ after imposing a declared 0.05 model-systematics floor. The
integer quadratic law is therefore retained rather than fitting the exponent.
Leave-one-subtype-out tests preserve the ordered sequence; the largest
measurement-interval discrepancy is the plusLRD held-out prediction, 0.923,
which is 0.012 above its measurement-only 84th percentile.

We use one soft 50~kK Cloudy basis and $\log U=-2.1$ for all sources. The
temperature denotes the incident spectral basis rather than a literal source
temperature. The closure predicts [O~III]$\lambda5007$/H$\beta$ of 0.95,
3.16, 4.43, and 5.07 and [O~III]/[O~II] of approximately 21.9 across the four
subtypes. Compared with the eight stacked ratios, the combined descriptive
$\chi^2$ is 5.77. All predicted He~II$\lambda4686$/H$\beta$ ratios lie between
0.0017 and 0.0154, below the independent population upper envelope of 0.1.
A 50--70~kK geometry-dependent hardness gradient worsens the same
[O~III] comparison to $\chi^2=17.0$. The data therefore favor a changing
dense--diffuse photon partition at nearly common soft incident hardness.

This comparison is a closure test, not a fully independent validation: the
stacked dense fractions were inferred from the same two [O~III] ratios. The
weak-He~II constraint is the independent hardness gate. Moreover, the
intrinsic predicted H$\alpha$/H$\beta$ ratios remain 2.77--3.15, so the large
observed Balmer decrements still require a separate attenuation or line-transfer
likelihood.

The subtype counts expose a distinct unresolved problem. With the present UV
flux gate, the predicted xLRD fraction is only 1\%, compared with approximately
12\% in the published stack sample. Replacing only that gate by an optical
5100-\AA{} gate raises the prediction to 37\%. Thus the red-population residual
is dominated by the uncalibrated multiband completeness mapping. We do not
interpret it as evidence for an additional formation route until the synthetic
photometry is propagated through survey masks and injection/recovery results.
